\section{Related Work}
\label{sec:related}

Tool-agent evaluation spans MCP benchmarks, API/tool-use benchmarks, task
generation, stable execution, and state- or trace-based scoring. We summarize
the main contrast here and give a detailed comparison in Appendix~\ref{app:rw}.
Prior work usually fixes at least one part of the setting: it ships a fixed
dataset, generates tasks from an imposed graph or plan, or scores an answer,
tool list, outcome state, or cached call trace. DynamicMCPBench targets the
complementary setting: a re-runnable framework over live or user-supplied MCP
servers, with tasks generated forward from successful executions and scored by
path-agnostic effects under deterministic replay.

Recent MCP benchmarks have expanded scale and realism: thousands of servers
and tools \citep{fan2025mcptoolbench++, mo2025livemcpbench, lei2025mcpverse},
live multi-step tasks \citep{wang2025mcp}, execution-grounded scoring
\citep{servers25mcp, gao2025mcp}, planted distractors \citep{bandi2026mcp},
stress tests \citep{wu2025mcpmark, yin2025livemcp, guo2026mcp}, GUI and
computer-use settings \citep{yan2025mcpworld, jia2025osworld}, and security
probes \citep{zhang2025mcp}. These benchmarks are valuable, but they remain
fixed task sets and usually score an answer, outcome, or tool choice. They do
not let practitioners re-run the benchmark on their own live servers, nor do
they score whether the required effects were produced along any valid path.

Older API/tool-use benchmarks study retrieval and selection
\citep{patil2024gorilla, qin2024toolllm}, confusable tools
\citep{huang2024metatool}, and call-structure correctness
\citep{patil2025berkeley}. Closest to us, \citet{yao2024tau} compare final
database state rather than exact paths and introduce pass$^k$ for reliability.
We keep this outcome-state intuition, but generalize it to path-agnostic
effects across many live MCP servers. We also ground every task in a real
successful trajectory, avoiding tool lists as unrecoverable prompt-level ground
truth \citep{qin2024toolllm}.

Task-generation work usually imposes structure and works backward: graph
construction, subgraph sampling, or generate-then-verify pipelines
\citep{shen2024taskbench, guo2026unitoolbench, liu2025mcpeval,
shi2025taskcraft}. Large real-server trajectories are also collected for
training \citep{xu2025toucan}. DynamicMCPBench instead explores first and
distills afterward: the successful trace becomes checkpoints, minefields, and a
partial order. To make live execution reproducible, we build on cached or
simulated tool environments
\citep{guo2024stabletoolbench, guo2025stabletoolbench, cheng2025travelbench}
and answer-agnostic state or trajectory scoring
\citep{lu2025toolsandbox, barres2025tau, kim2025beyond, zeng2026logigen,
chuang2026toward}.

Finally, tool surfaces and benchmark generation introduce confounds. MCP
descriptions are often low quality \citep{hasan2026model, wang2026docs}, many
exposed tools hurt selection \citep{gan2025rag}, and self-generated benchmarks
can inflate model scores \citep{yuan2026silencer}. These findings motivate our
headline setting: raw tool descriptions, multi-family authorship, and scoring
effects rather than final answers or curated tool lists. In short,
DynamicMCPBench is re-runnable, forward-generated, trace-grounded, and
effect-scored; \S\ref{sec:benchmark} details the framework, and
\S\ref{sec:results} reports the study.
